Let $e_i$ be the coordinate idempotents. A prime omits exactly one $e_i$, since $\sum e_i=1$ and $e_ie_j=0$ for $i\ne j$. It therefore has the form $A_1\times\cdots\times\mathfrak p_i\times\cdots\times A_n$. The prime correspondence for the projection identifies this set with $\operatorname{Spec}A_i$. It is $D(e_i)=V(1-e_i)$, hence open and closed.

Suppose $\operatorname{Spec}A=V(\mathfrak a)\sqcup V(\mathfrak b)$ with both parts nonempty. Then $\mathfrak a,\mathfrak b$ are proper, $\mathfrak a+\mathfrak b=A$, and $\mathfrak a\mathfrak b\subseteq\sqrt{(0)}$. Choose $a\in\mathfrak a$, $b\in\mathfrak b$ with $a+b=1$, and $n$ with $(ab)^n=0$. No maximal ideal contains both $a^n,b^n$, so $ua^n+vb^n=1$ for some $u,v\in A$. Then $e=ua^n$ satisfies $e(1-e)=0$. Also $e\ne0,1$, since $e\in\mathfrak a$ and $1-e\in\mathfrak b$.

Conversely, if $e^2=e\ne0,1$, the ideals $(e),(1-e)$ are comaximal and have zero intersection. The Chinese remainder theorem gives $A\cong A/(e)\times A/(1-e)$, with both factors nonzero. The first paragraph makes its spectrum disconnected. In particular, a local spectrum is connected by \exref{1.12}.
